
LLM self-repair has emerged as an approach to code generation that iteratively refines generated code using execution feedback until tests pass or a budget is exhausted. Prior work reports 10-17\% relative improvement on standard benchmarks through this mechanism. Recent studies have explored augmenting execution feedback with static analysis, demonstrating reductions in security and reliability issues.

Existing studies conflate feedback content with feedback presentation. When comparing static analysis feedback against execution feedback, prior work necessarily varies the information provided. This leaves a question unanswered: does the order in which feedback is presented independently affect repair quality, holding content constant?

This gap matters because presentation order is a free variable in any feedback-based repair system. If order affects quality, current systems may leave performance gains unexploited.

The hypothesis tested here is that static-first feedback ordering creates a coarse-to-fine repair trajectory. When static analysis errors appear before execution failures, the LLM resolves surface-level issues before tackling semantic errors. The key test is not that static-first provides different information, but that it structures the repair process differently.

The main contributions are:

\begin{enumerate}
\item A matched-content experimental design that isolates feedback ordering from information volume by ensuring byte-identical content across conditions.
\end{enumerate}

\begin{enumerate}
\item Evidence that static-to-execution ordering achieves 29.02\% relative pass@1 improvement over execution-to-static ordering (p<$10^{-5})$ on 664 problems.
\end{enumerate}

\begin{enumerate}
\item Mechanism analysis showing the effect operates through regression prevention (38\% reduction, p=0.0198) rather than early-gain acceleration (p=0.859).
\end{enumerate}

